
\documentclass{article}
\usepackage{colortbl}
\usepackage{makecell}
\usepackage{multirow}
\usepackage{supertabular}

\begin{document}

\newcounter{utterance}

\centering \large Interaction Transcript for game `codenames', experiment `opponent\_difficult', episode 4 with reflexion{-}Qwen3.5{-}4B.
\vspace{24pt}

{ \footnotesize  \setcounter{utterance}{1}
\setlength{\tabcolsep}{0pt}
\begin{supertabular}{c@{$\;$}|p{.15\linewidth}@{}p{.15\linewidth}p{.15\linewidth}p{.15\linewidth}p{.15\linewidth}p{.15\linewidth}}
    \# & $\;$A & \multicolumn{4}{c}{Game Master} & $\;\:$B\\
    \hline

    \theutterance \stepcounter{utterance}  
    & & & \multicolumn{2}{p{0.3\linewidth}}{
        \cellcolor[rgb]{0.9,0.9,0.9}{
            \makecell[{{p{\linewidth}}}]{
                \texttt{\tiny{[GM$|$GM]}}
                {'hidden': {'team': ['robin', 'pumpkin', 'unicorn', 'smuggler', 'grace'], 'innocent': ['figure', 'queen', 'bank', 'hood', 'circle'], 'opponent': [], 'assassin': ['net']}, 'revealed': {'team': {'team': [], 'innocent': [], 'opponent': [], 'assassin': []}, 'opponent': {'team': [], 'innocent': [], 'opponent': [], 'assassin': []}}}
            }
        }
    }
    & & \\ \\

    \theutterance \stepcounter{utterance}  
    & & \multicolumn{4}{p{0.6\linewidth}}{
        \cellcolor[rgb]{0.9,0.9,0.9}{
            \makecell[{{p{\linewidth}}}]{
                \texttt{\tiny{[P1$\langle$GM]}}
                \texttt{You have attempted this same task before and did not succeed. These are the plans you wrote after those attempts.} \\
\texttt{Trial \#1: New plan: The last attempt was aborted because the clue word "fly" was not semantically related to the target words "robin", "unicorn", or "spider" (as "robin" and "unicorn" are not typically associated with flying in the same direct way as "spider" is not a flyer, and "robin" is a bird but the clue must relate to the *team's* words collectively or individually without being a distractor). I will choose a single{-}word clue that is a strong, direct synonym or category for at least two of the target words, ensuring it is not a distractor or a kill word. For example, I will use "bird" to relate to "robin" and "note" (as in a bird's song or a musical note, though "note" is also a target, so I must avoid using a target word as the clue). A better choice is "song" for "robin" and "sound" (since a robin makes a sound), or "music" for "grace" and "note". Let's select "music" which relates to "grace" (grace notes) and "note".} \\
\\ 
\texttt{CLUE: music} \\
\texttt{TARGETS: grace, note} \\
\\ 
\texttt{Trial \#2: New plan: I failed because I grouped unrelated words (unicorn, spider, robin) under the clue "bird," causing my teammate to guess an opponent's word ("duck") instead of my targets. I must ensure every word in my TARGETS list is strongly and exclusively associated with my CLUE word, avoiding any ambiguity that leads to opponent or distractor guesses. For the remaining words (robin, pumpkin, unicorn, smuggler, grace), I will select a clue that connects at least two of them without overlap with the opponent's list. "Grace" and "robin" can be linked to "song" (grace notes, robin's song), or "pumpkin" and "robin" to "fall" (pumpkins are autumn, robins migrate in fall). I will choose "fall" for pumpkin and robin, or "song" for grace and robin, ensuring the clue is safe and covers multiple targets.} \\
\\ 
\texttt{CLUE: song} \\
\texttt{TARGETS: grace, robin} \\
\\ 
\\ 
\texttt{Provide one single{-}word clue that relates to some of your team's words from the list below. You can choose to how many and to which words your clue relates to.} \\
\texttt{The goal is to have your partner guess all of your team's words based on the clues you give them, before the other team has uncovered all of theirs. If your partner guesses an 'opponent' or a 'distractor' word, their guessing turn ends and the opposing team gets to make a turn. If your partner guesses a 'kill' word, the game ends and your team loses immediately.} \\
\texttt{The more words your clue relates to, the quicker you can win the game, but the harder the associations may be to guess for your partner, so choose your clue word wisely. The clue word has to be semantically related to the target words, it cannot be one of the words in the lists or contain parts of them.} \\
\\ 
\texttt{Always give your single{-}word clue and your comma{-}separated list of related target words in the following format and make your answers as short as possible, never include any other text than is required in this form:} \\
\\ 
\texttt{CLUE: <WORD>} \\
\texttt{TARGETS: <TARGETS>} \\
\\ 
\texttt{Your team words are: robin, pumpkin, unicorn, smuggler, grace.} \\
\texttt{Your opponent's team words are: .} \\
\texttt{Distractor words are: figure, queen, bank, hood, circle.} \\
\texttt{The kill words are: net.} \\
            }
        }
    }
    & \\ \\

    \theutterance \stepcounter{utterance}  
    & \multicolumn{4}{p{0.6\linewidth}}{
        \cellcolor[rgb]{0.9,0.9,0.9}{
            \makecell[{{p{\linewidth}}}]{
                \texttt{\tiny{[P1$\rangle$GM]}}
                \texttt{CLUE: song} \\
\texttt{TARGETS: grace, robin} \\
            }
        }
    }
    & & \\ \\

    \theutterance \stepcounter{utterance}  
    & & & \multicolumn{2}{p{0.3\linewidth}}{
        \cellcolor[rgb]{0.9,0.9,0.9}{
            \makecell[{{p{\linewidth}}}]{
                \texttt{\tiny{[GM$|$GM]}}
                \texttt{song} \\
            }
        }
    }
    & & \\ \\

    \theutterance \stepcounter{utterance}  
    & & & \multicolumn{2}{p{0.3\linewidth}}{
        \cellcolor[rgb]{0.9,0.9,0.9}{
            \makecell[{{p{\linewidth}}}]{
                \texttt{\tiny{[GM$|$GM]}}
                ['grace', 'robin']
            }
        }
    }
    & & \\ \\

    \theutterance \stepcounter{utterance}  
    & & & \multicolumn{2}{p{0.3\linewidth}}{
        \cellcolor[rgb]{0.9,0.9,0.9}{
            \makecell[{{p{\linewidth}}}]{
                \texttt{\tiny{[GM$|$GM]}}
                {'word': 'grace', 'assignment': 'team'}
            }
        }
    }
    & & \\ \\

    \theutterance \stepcounter{utterance}  
    & & & \multicolumn{2}{p{0.3\linewidth}}{
        \cellcolor[rgb]{0.9,0.9,0.9}{
            \makecell[{{p{\linewidth}}}]{
                \texttt{\tiny{[GM$|$GM]}}
                {'word': 'robin', 'assignment': 'team'}
            }
        }
    }
    & & \\ \\

    \theutterance \stepcounter{utterance}  
    & & & \multicolumn{2}{p{0.3\linewidth}}{
        \cellcolor[rgb]{0.9,0.9,0.9}{
            \makecell[{{p{\linewidth}}}]{
                \texttt{\tiny{[GM$|$GM]}}
                \texttt{opponent has won} \\
            }
        }
    }
    & & \\ \\

    \theutterance \stepcounter{utterance}  
    & \multicolumn{4}{p{0.6\linewidth}}{
        \cellcolor[rgb]{0.95,0.95,0.95}{
            \makecell[{{p{\linewidth}}}]{
                \texttt{\tiny{[P2$\rightarrow$P1]}}
                \texttt{Post{-}episode Reflexion (not sent to GM):} \\
\texttt{New plan: I failed because "song" is too ambiguous; while it relates to "robin" (song) and "grace" (grace notes), it could also be associated with "unicorn" (mythical creatures often depicted singing) or "smuggler" (smuggling songs), potentially leading my partner to guess a distractor or an opponent's word. I must select a clue that is uniquely and strongly associated with at least two of my specific target words without any overlap with the distractor list (figure, queen, bank, hood, circle) or the kill word (net). I will choose "fall" as the clue, which strongly and exclusively relates to "pumpkin" (autumn harvest) and "robin" (migratory bird), ensuring my partner guesses only these two targets.} \\
\\ 
\texttt{CLUE: fall} \\
\texttt{TARGETS: pumpkin, robin} \\
            }
        }
    }
    & & \\ \\

\end{supertabular}
}

\end{document}
